The present study has four deliberate limitations. First, the verifier is a released RF model; no extrapolation candidate has been confirmed by molecular dynamics or experiment. Second, the data source is a single HTP-MD protocol, so the conclusions are conditional on its chemistry and simulation support. Third, structural-OOD evaluation is blocked because scaffold/family annotations and matched continuous/binary checkpoints are not yet available. Fourth, the six-property benchmark is a prediction-phase audit and should not be interpreted as a validated multi-property physical model.

The next study should therefore begin with provenance-preserving expansion of the labelled conductivity corpus, followed by a representation ablation and a closed-loop acquisition experiment. A larger unlabelled database such as OPoly26 may be useful for geometry or force-field pretraining, but it cannot directly expand the conductivity label space. Active learning should be evaluated after MD calibration and with a pre-registered acquisition criterion.\

\placeholder{Insert final MD validation results, external-data calibration, structural-OOD split and active-learning results when available.}
